\documentclass[11pt,a4paper]{article}
\usepackage[margin=2.4cm]{geometry}
\usepackage{amsmath,amsthm,mathtools}
\usepackage{fontspec}
\usepackage{unicode-math}
\setmainfont{STIX2Text}[Path=fonts/, Extension=.otf,
  UprightFont=*-Regular, ItalicFont=*-Italic,
  BoldFont=*-Bold, BoldItalicFont=*-BoldItalic]
\setmathfont{STIX2Math}[Path=fonts/, Extension=.otf]
\usepackage[protrusion=true,expansion=false]{microtype}
\usepackage{booktabs,enumitem}
\usepackage{tikz}
\usetikzlibrary{arrows.meta,positioning,calc,fit,backgrounds}
\definecolor{tierspend}{RGB}{200,60,60}
\definecolor{tierfull}{RGB}{210,140,40}
\definecolor{tierin}{RGB}{60,130,180}
\definecolor{tierrecv}{RGB}{70,150,90}
\usepackage[psdextra,hidelinks]{hyperref}

\emergencystretch=3em
\hbadness=10000
\hfuzz=2pt

\newtheorem{theorem}{Theorem}[section]
\newtheorem{lemma}[theorem]{Lemma}
\newtheorem{proposition}[theorem]{Proposition}
\newtheorem{corollary}[theorem]{Corollary}
\theoremstyle{definition}
\newtheorem{definition}[theorem]{Definition}
\newtheorem{construction}[theorem]{Construction}
\newtheorem{assumption}[theorem]{Assumption}
\newtheorem{remark}[theorem]{Remark}
\newtheorem{example}[theorem]{Example}

% kind-coded theorem blocks, palette shared with the figures and the website:
% definitions/constructions blue (tierin), assumptions amber (tierfull),
% proved statements a neutral bar; remarks and examples run unframed
\usepackage{mdframed}
\providecolor{tierin}{RGB}{60,130,180}
\providecolor{tierfull}{RGB}{210,140,40}
\mdfdefinestyle{kindbar}{topline=false,bottomline=false,rightline=false,
  linewidth=1.5pt,innerleftmargin=9pt,innerrightmargin=4pt,
  innertopmargin=5pt,innerbottommargin=6pt,leftmargin=0pt,rightmargin=0pt,
  skipabove=10pt,skipbelow=10pt}
\surroundwithmdframed[style=kindbar,linecolor=tierin!70,backgroundcolor=tierin!4]{definition}
\surroundwithmdframed[style=kindbar,linecolor=tierin!70,backgroundcolor=tierin!4]{construction}
\surroundwithmdframed[style=kindbar,linecolor=tierfull!80,backgroundcolor=tierfull!5]{assumption}
\surroundwithmdframed[style=kindbar,linecolor=black!30]{theorem}
\surroundwithmdframed[style=kindbar,linecolor=black!30]{lemma}
\surroundwithmdframed[style=kindbar,linecolor=black!30]{proposition}
\surroundwithmdframed[style=kindbar,linecolor=black!30]{corollary}

\usepackage{fancyhdr}
\pagestyle{fancy}
\fancyhf{}
\fancyhead[L]{\footnotesize\scshape The Zcash Arboretum}
\fancyhead[R]{\footnotesize\itshape\nouppercase{\leftmark}}
\fancyfoot[C]{\footnotesize\thepage}
\renewcommand{\headrulewidth}{0.3pt}
\renewcommand{\sectionmark}[1]{\markboth{\thesection\ \ #1}{}}

\newcommand{\F}{\mathbb{F}}
\newcommand{\Z}{\mathbb{Z}}
\newcommand{\N}{\mathbb{N}}
\newcommand{\G}{\mathbb{G}}
\newcommand{\E}{\mathbb{E}}
\newcommand{\PP}{\mathbb{P}}
\renewcommand{\Pr}{\mathrm{Pr}}
\newcommand{\Fp}{\F_p}
\newcommand{\Fq}{\F_q}
\newcommand{\ip}[2]{\langle #1,\, #2 \rangle}
\newcommand{\Comm}{\mathsf{Commit}}
\newcommand{\Open}{\mathsf{Open}}
\newcommand{\Setup}{\mathsf{Setup}}
\newcommand{\Verify}{\mathsf{Verify}}
\newcommand{\negl}{\mathsf{negl}}
\newcommand{\poly}{\mathsf{poly}}
\newcommand{\PPT}{\mathsf{PPT}}
\newcommand{\Adv}{\mathcal{A}}
\newcommand{\Prov}{\mathcal{P}}
\newcommand{\Ver}{\mathcal{V}}
\newcommand{\Sim}{\mathcal{S}}
\newcommand{\Ext}{\mathcal{E}}
\newcommand{\Rel}{\mathcal{R}}
\newcommand{\Lang}{\mathcal{L}}
\newcommand{\nfsym}{\mathsf{nf}}
\newcommand{\cmsym}{\mathsf{cm}}
\newcommand{\cmx}{\mathsf{cmx}}
\newcommand{\cvsym}{\mathsf{cv}}
\newcommand{\rk}{\mathsf{rk}}
\newcommand{\Pallas}{\ensuremath{\mathsf{Pallas}}}
\newcommand{\Vesta}{\ensuremath{\mathsf{Vesta}}}
\newcommand{\Ep}{\ensuremath{\mathcal{E}}}
\newcommand{\NoteCommit}{\ensuremath{\mathsf{NoteCommit}}}
\newcommand{\Sinsemilla}{\ensuremath{\mathsf{Sinsemilla}}}
\newcommand{\ShortCommit}{\ensuremath{\mathsf{ShortCommit}}}
\newcommand{\Extract}{\ensuremath{\mathsf{Extract}}}
\newcommand{\Acc}{\ensuremath{\mathsf{Acc}}}
\newcommand{\GroupHash}{\ensuremath{\mathsf{GroupHash}}}
\newcommand{\ZH}{Z_H}
\newcommand{\ask}{\ensuremath{\mathsf{ask}}}
\newcommand{\aksym}{\ensuremath{\mathsf{ak}}}
\newcommand{\nksym}{\ensuremath{\mathsf{nk}}}
\newcommand{\fvk}{\ensuremath{\mathsf{fvk}}}
\newcommand{\ovk}{\ensuremath{\mathsf{ovk}}}
\newcommand{\dk}{\ensuremath{\mathsf{dk}}}
\newcommand{\pkd}{\ensuremath{\mathsf{pk_d}}}
\newcommand{\gd}{\ensuremath{\mathsf{g_d}}}
\newcommand{\rcv}{\ensuremath{\mathsf{rcv}}}
\newcommand{\rcm}{\ensuremath{\mathsf{rcm}}}
\newcommand{\rsk}{\ensuremath{\mathsf{rsk}}}

\renewenvironment{abstract}{%
  \small\begin{center}{\bfseries\abstractname}\end{center}\medskip\noindent\ignorespaces
}{\par\medskip}

\title{\textbf{\Huge Halo 2 Guide}\\[6pt]\large From a circuit to a checked proof}
\author{m@rek.onl}
\date{}
\begin{document}
\maketitle
\thispagestyle{empty}
\begin{abstract}
This volume of \emph{The Zcash Arboretum} constructs the Halo~2 proof system
used to prove shielded transaction relations.  It assumes the \emph{Math
Guide}'s fields, polynomials, linear algebra and Pasta curves, and the
\emph{Crypto Guide}'s commitments, security contracts and interactive
proofs.  We construct a way to convince a verifier that private data
satisfies public rules without disclosing that data.  The verifier's
complete calculation makes explicit what it checks, what the application
must supply, and which security conclusions require more than the
algebraic identities.
\end{abstract}
\tableofcontents
\newpage

\section{From a relation to a circuit}\label{sec:arithmetisation}

\subsection{The statement being proved}

The mathematical prerequisites are the \emph{Math Guide}, ``Prime fields
and modular arithmetic'', ``Linear algebra and uniform masks'',
``Polynomials and evaluation'', ``Evaluation domains and fast transforms'',
and ``Pasta curves and their encodings''.  We use their fields, vectors,
polynomials and curve groups with the same definitions.

A proof establishes a precisely selected relation, not an informal claim
that a transaction is good.  Let $u$ be a public input, $w$ a private
witness, and $\Rel(u,w)$ a deterministic predicate.
The prover receives both inputs; the verifier receives $u$ and a proof.
Completeness asks that an honestly constructed proof be accepted when
$\Rel(u,w)=1$.  Knowledge soundness asks for a witness-extraction guarantee
against an efficient prover that succeeds sufficiently often.  Zero
knowledge asks that the verifier's view be simulatable from the public
input alone.  These distinct contracts are defined in the \emph{Crypto
Guide}, ``Security contracts and reductions'' and ``Knowledge, simulation,
and Schnorr''.

The application chooses $\Rel$.  Halo~2 represents that relation by an
\emph{arithmetic circuit}: field-valued cells connected by polynomial
equations which must hold simultaneously.  Its soundness cannot repair
an omitted bound on a value, a
wrong hash input, or a public value left unconnected to the witness.
There are therefore two obligations:
\[
\begin{aligned}
 \text{circuit satisfaction}&\Longrightarrow\Rel(u,w)=1,\\
 \text{accepted proof}&\Longrightarrow
 \text{the proof system's knowledge guarantee}.
\end{aligned}
\]
The first is about the circuit.  The second is about the argument system.
Neither implication follows merely from the other.

We use the Zcash Halo~2 implementation, specifically
\texttt{halo2\_proofs} at commit
\begin{center}
\texttt{261faaccd5a30c19bb8468600841a0dac305adf5}.
\end{center}
The pinned source is the authority for the folding and transcript
conventions below.  Other systems named Halo~2 can use different polynomial
commitments or transcripts.  The implementation sources are linked in
\S\ref{sec:halo-sources}.  This construction does not verify earlier
proofs inside another proof.

\subsection{The field and the table}

Let $\F$ be the field in which the circuit computes.  For the relevant
Pasta instantiation it is the Pallas base field, equivalently the Vesta
scalar field.  The commitment group used below is Vesta, not Pallas.
The \emph{Math Guide}, ``Pasta curves and their encodings'', gives both
fields and their encodings.

A circuit is a fixed arrangement of columns and rows.  An \emph{advice}
column holds witness-dependent values.  A \emph{fixed} column holds
circuit constants.  An \emph{instance} column holds public inputs, padded
with zeroes to the circuit size.  A \emph{gate} is a polynomial expression
in selected cells which must equal zero.  A query can select the current
row or a fixed rotation, such as the next row.  Rows are indexed modulo
the row count; a next-row query at the last row wraps to row zero.

A \emph{verification key} is public data fixing this circuit's columns,
equations, constants and wiring.  The verifier must use the intended key,
not a different circuit selected by the prover.  A fixed selector is a
column which turns a gate on at its intended rows.
For example,
\[
 q_{\mathrm{mul}}(a b-c)=0
\]
imposes multiplication where the selector is one and imposes nothing
where it is zero.  In the implementation, selectors may be compressed
into fixed columns and replaced by polynomial expressions.  The
verification key fixes the resulting expressions; they are not choices
the prover can change after seeing a challenge.

Assignments are not constraints.  Computing a value in the witness
generator only tells the honest prover what to put in a cell.  A
malicious prover need not run that generator.  Two cells assigned the
same host-language variable must still be joined by an equality
constraint if their equality matters.  Similarly, a witness calculated
from a public input must be constrained to that input.

\subsection{A complete small computation}

Consider the relation
\[
 \Rel(t,s)=1\quad\Longleftrightarrow\quad s^3+s+5=t
 \quad\text{in }\F_{97}.
\]
The public input is $t=35$; one witness is $s=3$.  We use four rows and
the gate
\begin{equation}\label{eq:toy-gate}
 q_Mab+q_La+q_Rb+q_Oc+q_C+\pi=0.
\end{equation}
The column $\pi$ is public.  All five $q$ columns are fixed.
\[
\begin{array}{c|rrr|rrrrr|r}
 i&a&b&c&q_M&q_L&q_R&q_O&q_C&\pi\\ \hline
 0&3&3&9&1&0&0&-1&0&0\\
 1&9&3&27&1&0&0&-1&0&0\\
 2&27&3&30&0&1&1&-1&0&0\\
 3&30&0&0&0&1&0&0&5&-35
\end{array}
\]
The negative entries denote field elements: for example $-1=96$.
The equalities are
\[
 a_0=b_0=b_1=b_2,\qquad
 c_0=a_1,\qquad c_1=a_2,\qquad c_2=a_3.
\]
These make the four local equations one computation:
$s^2$, then $s^3$, then $s^3+s$, then comparison with $t-5$.
Without the copy equalities, the rows could describe unrelated numbers.

This four-row table isolates the arithmetic.  It does not reserve
blinding rows---unused rows filled with randomness to conceal witness
values---or instantiate a secure curve, so it is not a production
Halo~2 proof.  Its purpose is to make the polynomial calculations
independently reproducible.  The verifier and masking rules for an actual
circuit are introduced below.

\subsection{Integers inside a field}

The circuit sees residues, not unbounded integers.  A field equality
$a+b=c$ also holds when the intended integer sum differs by the field
modulus.  To interpret a field element as a $k$-bit nonnegative integer,
one can constrain bits $b_i$ by
\[
 b_i(b_i-1)=0,\qquad
 a=\sum_{i=0}^{k-1}2^i b_i,
\]
and require $2^k$ not to exceed the modulus.  A \emph{lookup} checks that
a value belongs to a specified table.  Such table-membership checks can
replace individual bit constraints when enforcing a range of integers.
Bounds on intermediate sums
must also rule out wraparound.  The \emph{Math Guide}, ``Language,
integers, and encodings'' and ``Prime fields and modular arithmetic'',
proves the required integer-to-field statements.  This obligation is
essential whenever a circuit represents money or encoded byte strings.

\subsection{The commitment interface}

The \emph{Crypto Guide}, ``Groups, commitments, and Sinsemilla'',
constructs commitments to vectors of scalars.  Committing to a
polynomial of degree less than $n$ means applying that construction to
its coefficient vector, padded to length $n$:
\[
 C=\sum_{i=0}^{n-1}[a_i]G_i+[r]W
 \quad\text{for }p(X)=\sum_{i=0}^{n-1}a_iX^i.
\]
The public points $G_i,W$ belong to a group with scalar field $\F$.
The scalar $r$ is the commitment blind.  The lower volume's hiding and
binding contracts apply to these vectors.  The concrete parameter
derivation and the algorithm for proving an evaluation are constructed
in \S\ref{sec:ipa}; until then, the arithmetic reductions describe the
identities that such evaluations must authenticate.

The \emph{Crypto Guide}, ``Knowledge, simulation, and Schnorr'', also
introduces the Fiat--Shamir transformation: deterministic hash-derived
challenges replace an interactive verifier's random challenges.
We first explain the arithmetic with random challenges.  A
\emph{transcript} is the ordered record of public inputs, messages and
challenges.  Its concrete byte construction is given in
\S\ref{sec:transcript}; its security is not assumed to follow from a
single random-point calculation.

\section{Columns as polynomials}\label{sec:polynomials}

\subsection{A row is an evaluation point}

Choose $n=2^k$ dividing $|\F|-1$ and a primitive $n$th root of unity
$\omega$.  The evaluation domain is
\[
 H=\{1,\omega,\ldots,\omega^{n-1}\},\qquad Z_H(X)=X^n-1.
\]
For each column with entries $a_i$, there is exactly one polynomial
$a(X)$ of degree less than $n$ satisfying $a(\omega^i)=a_i$.
The \emph{Math Guide}, ``Polynomials and evaluation'', constructs it as
$\sum_i a_iL_i(X)$, where
\[
 L_i(X)=\frac{\omega^i}{n}\frac{X^n-1}{X-\omega^i}.
\]
The quotient is a polynomial; the displayed rational expression is
evaluated by cancellation at a domain point.  A rotation by $r$ rows is
the polynomial $a(\omega^rX)$.  It has the same degree as $a$.

Substituting the column polynomials into a gate produces a polynomial
$g(X)$.  Its row constraints hold exactly when $g$ vanishes on $H$.
Since the elements of $H$ are distinct,
\begin{equation}\label{eq:divisibility}
 \forall h\in H,\ g(h)=0
 \quad\Longleftrightarrow\quad
 \exists q(X):\ g(X)=Z_H(X)q(X).
\end{equation}
This is polynomial divisibility, not pointwise division on $H$, where
the denominator would be zero.

\subsection{The four-row table in coefficient form}

In the toy example, choose $\omega=22$.  Its successive powers modulo
$97$ are $1,22,96,75$.  Interpolation gives
\[
\begin{aligned}
 a(X)&=90+61X+22X^2+24X^3,\\
 b(X)&=75+32X+25X^2+65X^3,\\
 c(X)&=65+16X+3X^2+22X^3.
\end{aligned}
\]
Interpolate the fixed and public columns in exactly the same way and
substitute them into \eqref{eq:toy-gate}.  The resulting gate polynomial
is
\[
 g(X)=(X^4-1)
 (61+70X+43X^2+47X^3+81X^4+46X^5).
\]
Every coefficient and the following evaluations are checked by
\texttt{research/check\_foundations.py}.  At the off-domain point $x=20$,
\[
 Z_H(20)=46,\qquad q(20)=67,\qquad g(20)=75,
 \qquad 46\cdot67=75\pmod{97}.
\]

Now change only $c_0$ from $9$ to $10$.  Division of the new gate
polynomial by $X^4-1$ leaves remainder
\[
 r(X)=24(1+X+X^2+X^3).
\]
Its roots in $\F_{97}$ are $22,75,96$: exactly the three unchanged rows.
At $x=20$, the altered numerator is $20$, while its Euclidean quotient
times $Z_H(20)$ is $64$.  The test rejects.  The changed table also breaks
a copy equality, but this calculation concerns the gate check alone.

For this particular remainder, a uniform field point is a root with
probability $3/97$.  That is not a security estimate for Halo~2.  The
field is deliberately tiny, the commitment argument has not yet been
applied, and the production verifier has exceptional cases such as
$x\in H$ which must be accounted for separately.

\subsection{Combining identities}

Suppose the circuit and its auxiliary arguments require
$g_0,\ldots,g_{M-1}$ to vanish on $H$.  After their defining commitments
are fixed, sample a challenge $y$ and combine them in Horner order:
\[
 N_y(X)=
 y^{M-1}g_0(X)+y^{M-2}g_1(X)+\cdots+g_{M-1}(X).
\]
The order is the verification key's gate order, followed by permutation
and lookup constraints, repeated for each circuit instance in a joint
proof.

If one row violates any of these identities, at that row the combined
expression is a nonzero polynomial in $y$ of degree at most $M-1$.
For fixed constraints and an independent uniform challenge, accidental
cancellation has probability at most $(M-1)/|\F|$, by the root bound in
the \emph{Math Guide}, ``Polynomials and evaluation''.  This is a local
interactive argument.  An actual Fiat--Shamir security bound must also
account for how the adversary obtains and retries challenges.

After receiving $y$, the prover commits to a quotient $h$ intended to
satisfy
\[
 N_y(X)=(X^n-1)h(X).
\]
The verifier subsequently samples $x$ and checks the equality at $x$.
If the difference polynomial is nonzero and has degree at most $D$,
an independent uniform $x$ hides that difference with probability at
most $D/|\F|$.  The timing matters: choosing $h$ after learning $x$ would
allow a prover to fit one point without satisfying the identity.

\subsection{Degree, cosets and quotient chunks}

Let $d$ be the configured constraint-degree bound.  It is the maximum
of the permutation and lookup requirements, all gate degrees, and any
explicitly configured minimum degree.  In counting expression degrees,
every queried column counts as degree one; selectors and boundary
polynomials count too.  Thus $d$ need not be attained by an expression.
Substituting degree-less-than-$n$ column polynomials gives
\[
 \deg N_y\le d(n-1).
\]
For a satisfying witness, division by the degree-$n$ polynomial $Z_H$
therefore gives $\deg h<(d-1)n$.  Write
\begin{equation}\label{eq:quotient-chunks}
 h(X)=\sum_{j=0}^{d-2}X^{jn}h_j(X),\qquad \deg h_j<n.
\end{equation}
Each chunk now fits the same commitment parameters as a column.
There are always $d-1$ chunk commitments, including any zero high
chunks; proof shape does not reveal the actual degree.

The prover evaluates the constraints on an extended coset of size
\[
 N=2^{\lceil\log_2((d-1)n)\rceil}.
\]
The implementation uses $\zeta\langle\omega_{\rm ext}\rangle$, where
$\zeta$ has order three and $\omega_{\rm ext}$ has order $N$.
The coset is disjoint from $H$, since a nontrivial element of order three
cannot belong to a group of power-of-two order.  Pointwise division by
$X^n-1$ is therefore defined there.  An inverse transform recovers the
quotient coefficients.  The \emph{Math Guide}, ``Evaluation domains and
fast transforms'', supplies these transforms.

The extended size accommodates the \emph{quotient}, not necessarily the
numerator.  For $d=9$, it is $8n$, although the numerator can have degree
close to $9n$.  Evaluating the constraint expression directly on the
coset and dividing is sufficient for an honest witness.  For an invalid
witness, the interpolated ratios need not be its Euclidean quotient;
acceptance still requires the independently challenged identity above.

Once $x$ is known, the verifier forms the commitment to
\[
 h_x(X)=\sum_{j=0}^{d-2}(x^n)^j h_j(X).
\]
The coefficients $(x^n)^j$ are now scalars.  This new polynomial has degree
less than $n$, and $h_x(x)=h(x)$.  Its expected value is
$N_y(x)/(x^n-1)$, calculated from the other claimed evaluations.  The proof
does not separately transmit each $h_j(x)$.

\subsection{Reserved rows}

Let $b$ be the number of blinding factors, and put $\ell=n-b-1$.
Rows $0,\ldots,\ell-1$ are usable.  Row $\ell$ holds the final value of
an auxiliary running product.  Rows $\ell+1,\ldots,n-1$ are product
blinding rows.  Define
\[
 L_{\rm blind}(X)=\sum_{i=\ell+1}^{n-1}L_i(X),
 \qquad A_{\rm act}(X)=1-L_\ell(X)-L_{\rm blind}(X).
\]
On $H$, the active-row polynomial is one exactly at usable rows.

Advice columns are randomised at \emph{all} $b+1$ unusable rows, including
row $\ell$.  Product polynomials retain their constrained value at row
$\ell$ and randomise only the last $b$ rows.  Gates must be disabled where
they would impose unintended constraints on these random cells.  A
rotation must likewise not make an active gate depend on a random
reserved cell unless that dependence is intended.

\section{The permutation argument}\label{sec:permutation}

\subsection{Copy constraints as a permutation}

Copy equalities partition selected cells into classes whose values must
agree.  Their endpoints must lie in the usable rows $0,\ldots,\ell-1$;
key generation rejects wiring to any other row.  Number the
equality-enabled columns $j=0,\ldots,m-1$.  Label the
cell in row $i$, column $j$, by
\[
 \operatorname{id}_{j,i}=\delta^j\omega^i.
\]
Choose $\delta$ so that these column cosets are disjoint.  Halo~2 uses
the field's designated $\mathsf{DELTA}$; the circuit must fit the supply
of disjoint cosets.  Distinct labels identify distinct cells even when
their witness values happen to coincide.

For each equality class, choose one cycle through its cells.  Cells
outside a nontrivial class are fixed points.  Together these cycles
define a permutation $\sigma$.  The fixed polynomial $s_j$ contains the
label of the successor:
\[
 s_j(\omega^i)=\operatorname{id}_{\sigma(j,i)}.
\]
Both the wiring and the commitments to these polynomials belong to the
verification key.

\begin{lemma}[Why a labelled product checks equality]
Let $v_c$ be the value at cell $c$, and let $\lambda_c$ be its distinct
label.  The formal identity
\[
 \prod_c(v_c+\beta\lambda_c+\gamma)
 =
 \prod_c(v_c+\beta\lambda_{\sigma(c)}+\gamma)
\]
holds as a polynomial in $\beta,\gamma$ if and only if values are
constant on every cycle of $\sigma$.
\end{lemma}
\begin{proof}
Fix a cell $c$ and substitute $\gamma=-v_c-\beta\lambda_c$.
The left product is zero.  The right becomes a product of polynomials
in $\beta$, so one of its factors must be the zero polynomial: the
product of nonzero polynomials over a field is nonzero.  Its coefficient
of $\beta$ requires $\lambda_{\sigma(d)}=\lambda_c$ for that factor's
cell $d$, and its constant term requires $v_d=v_c$.  Distinct labels
therefore give $d=\sigma^{-1}(c)$.  Doing this for every $c$ proves
constancy on each cycle.  Conversely, these equalities merely reorder
the factors.  The polynomial facts used here are proved in the
\emph{Math Guide}, ``Polynomials and evaluation''.
\end{proof}

For actual random field values $\beta,\gamma$, a nonzero difference
can vanish accidentally.  The polynomial root bound controls that
event; the product identity is not an unconditional test at one
challenge pair.  The advice commitments precede these challenges.

\subsection{The running product}

For one set of equality columns, an honest prover starts with $Z(1)=1$
and, for usable rows, computes
\[
 Z(\omega^{i+1})=
 Z(\omega^i)
 \frac{\prod_j(a_j(\omega^i)+\beta\delta^j\omega^i+\gamma)}
      {\prod_j(a_j(\omega^i)+\beta s_j(\omega^i)+\gamma)}.
\]
This formula describes the ordinary case in which all denominator
factors are nonzero.  It telescopes: the final product equals one
precisely when the two labelled products agree.

The verifier does not receive all these intermediate values.  It
checks a polynomial recurrence, multiplied by $A_{\rm act}$:
\begin{equation}\label{eq:permutation-recurrence}
\begin{split}
 A_{\rm act}(X)\biggl[
  Z(\omega X)\prod_j(a_j(X)+\beta s_j(X)+\gamma)\\
  {}-Z(X)\prod_j(a_j(X)+\beta\delta^jX+\gamma)
 \biggr]=0.
\end{split}
\end{equation}
Boundary constraints fix the initial and final values.  These
expressions join the combined numerator $N_y$; checking them uses the
same quotient machinery as checking an arithmetic gate.

\subsection{Chunking the columns}

A product over $m$ columns would make
\eqref{eq:permutation-recurrence} have expression degree $m+2$:
one for the active selector, one for $Z$, and $m$ factors.
For a circuit degree $d$, the implementation therefore puts at most
$d-2$ columns in a chunk.  Let $J_t$ be the columns of chunk $t$, and
let $Z_t$ be its product polynomial.

The first chunk starts at one.  Each subsequent chunk starts at the
previous chunk's final usable product.  The exact boundary constraints
are
\begin{align}
 L_0(X)(1-Z_0(X))&=0,\label{eq:perm-first}\\
 L_0(X)\bigl(Z_t(X)-Z_{t-1}(\omega^\ell X)\bigr)&=0
       &&(t>0),\label{eq:perm-link}\\
 L_\ell(X)\bigl(Z_{T-1}(X)^2-Z_{T-1}(X)\bigr)&=0.
       \label{eq:perm-last}
\end{align}
Each chunk also has recurrence \eqref{eq:permutation-recurrence} with
$j\in J_t$.  Here $T$ is the number of chunks.  At row zero,
$\omega^\ell X$ addresses row $\ell$, which explains the rotated
query in \eqref{eq:perm-link}.

The final constraint is \emph{Boolean}: it permits zero or one.  It is
not written as $Z_{T-1}=1$.  Off the exceptional set where a numerator or
denominator factor vanishes, every running product is nonzero, so this
Boolean constraint forces one.  A soundness argument must include the
exceptional event rather than silently replace the implemented
constraint by a stronger one.

The verifier needs each $Z_t$ at $x$ and $\omega x$, and all but the last
at $\omega^\ell x$.  The latter is also the rotation $-(b+1)$, because
$\omega^n=1$.  It needs every fixed $s_j$ at $x$.
At degree nine, fifteen equality-enabled columns require three
seven-column-capacity chunks, not two.  Columns can be advice, instance
or fixed; counting advice columns alone is insufficient.

\section{The lookup argument}\label{sec:lookups}

\subsection{Membership, not equal multiplicity}

A lookup requires a row's tuple to occur among the rows of a table.
For example, a table containing $0,\ldots,255$ constrains a value to
one byte without eight separate bit equations.  A table row can be
used repeatedly.  A \emph{multiset} records values together with their
multiplicities, meaning how many times each occurs.  The input multiset
therefore need not have the same multiplicities as the table.

Let the input tuple be $(a_0,\ldots,a_{r-1})$ and the table tuple
$(s_0,\ldots,s_{r-1})$.  After committing to the advice columns, sample
$\theta$ and compress each tuple in Horner order:
\[
 A=\sum_{j=0}^{r-1}\theta^{r-1-j}a_j,\qquad
 S=\sum_{j=0}^{r-1}\theta^{r-1-j}s_j.
\]
A pair of distinct, fixed tuples collides under this compression with
probability at most $(r-1)/|\F|$.  To protect all relevant pairs, the
analysis must count them and combine their bounds.  The single-pair
bound is not the error of the whole lookup argument.

Lookup expressions exist at every usable row.  Disabling a conditional
lookup must map the inactive input to a tuple actually present in the
table, often an all-zero tuple.  A selector which merely sets a value
to zero is not enough if zero is absent from the table.

\subsection{Arranging equal runs}

The prover permutes the compressed input values into $A'$ with equal
values adjacent.  It permutes the table values into $S'$ so that the
first row of each equal run in $A'$ matches the table value on that row.
Unused table entries fill the remaining positions.

For example, with input and table values
\[
 A=(3,1,3,3),\qquad S=(1,2,3,4),
\]
a valid arrangement is
\[
 A'=(1,3,3,3),\qquad S'=(1,3,2,4).
\]
The value three appears in three input rows but only one table row.
At the first three, the table matches.  At the next two, equality to
the preceding input suffices.

The conditions are
\[
 A'_0=S'_0,\qquad
 (A'_i-S'_i)(A'_i-A'_{i-1})=0\quad(0\le i<\ell).
\]
At row zero the separate first-row equation supplies the match,
regardless of the wrapped previous-row value.  At a later row, either
the current table entry matches or the preceding input is the same.
Induction from row zero shows that every input belongs to the table,
provided $A'$ and $S'$ really are permutations of their originals.
Sorting is an honest-prover construction; the verifier does not need
to prove a numerical ordering.

\subsection{Proving the two permutations}

The prover commits to $A'$ and $S'$ after $\theta$ but before the
independent challenges $\beta,\gamma$.  It forms a product $Z$ whose
ordinary-case recurrence is
\[
 Z(\omega^{i+1})=
 Z(\omega^i)
 \frac{(A(\omega^i)+\beta)(S(\omega^i)+\gamma)}
      {(A'(\omega^i)+\beta)(S'(\omega^i)+\gamma)}.
\]
The two challenges keep the input and table multiset tests separate.
As a formal polynomial identity, equality of the products in
$\beta,\gamma$ implies equality of each respective multiset: compare
the coefficient of the highest power of $\gamma$ to isolate the input
product, and the highest power of $\beta$ to isolate the table product.
The one-variable fact follows by matching the monic linear factors.

The exact constraint list is
\begin{align}
 L_0(1-Z)&=0,\label{eq:lookup-first}\\
 L_\ell(Z^2-Z)&=0,\label{eq:lookup-last}\\
 A_{\rm act}\bigl[
 Z(\omega X)(A'+\beta)(S'+\gamma)
 -Z(X)(A+\beta)(S+\gamma)\bigr]&=0,\label{eq:lookup-product}\\
 L_0(A'-S')&=0,\label{eq:lookup-run-first}\\
 A_{\rm act}(A'-S')\bigl(A'-A'(\omega^{-1}X)\bigr)&=0.
 \label{eq:lookup-runs}
\end{align}
Unmarked polynomials in these equations are evaluated at $X$.
As in the permutation argument, the Boolean endpoint forces one only
off the zero-factor exceptional set.

The prover randomises $A'$ and $S'$ in the last $b+1$ rows.  It
randomises $Z$ in the last $b$ rows while retaining its endpoint at
row $\ell$.  The verifier needs five evaluations per lookup:
\[
 Z(x),\quad Z(\omega x),\quad
 A'(x),\quad A'(\omega^{-1}x),\quad S'(x).
\]
Together with evaluations of the original input and table expressions,
these suffice to compute all five constraints at $x$.  Authentication
of those evaluations is the remaining task.

\section{Polynomial commitments and the masked inner-product argument}
\label{sec:ipa}

\subsection{Committing to a bounded-degree polynomial}

Let $\G$ be a prime-order group whose scalar field is $\F$.  Its identity
is $\mathcal O$.  Choose public generators
$G_0,\ldots,G_{n-1},W,U\in\G$ without known efficiently computable discrete
logarithm relations.  For $p(X)=\sum_{i=0}^{n-1}a_iX^i$ and blind $r\in\F$,
define
\[
 \Comm(p;r)=\sum_{i=0}^{n-1}[a_i]G_i+[r]W.
\]
This is the vector commitment from the \emph{Crypto Guide}, ``Groups,
commitments, and Sinsemilla''.  We use
$\ip{\boldsymbol a}{\boldsymbol G}=\sum_i[a_i]G_i$ for a scalar--point
inner product, and $\ip{\boldsymbol a}{\boldsymbol b}=\sum_i a_ib_i$
for a scalar--scalar inner product.

A uniform independent $r$ makes the commitment perfectly hiding:
multiplication by nonidentity $W$ is a bijection from $\F$ to $\G$.
Binding is computational, not information-theoretic.  The map from
$n+1$ scalars to one group element necessarily has many collisions;
the assumption is that a bounded adversary cannot find the relevant
distinct openings.  Binding alone does not supply an extractor for a
prover that sends group elements without their representations.

Halo~2 derives its parameters deterministically.  With the Pasta
hash-to-curve construction from the \emph{Crypto Guide}, ``Hash functions
and key derivation'', using domain \texttt{Halo2-Parameters}, it hashes
\[
\begin{array}{c|l}
 \text{point}&\text{message bytes}\\ \hline
 G_i&\mathtt{0x00}\parallel\operatorname{LE32}_{\rm bytes}(i)\\
 W&\mathtt{0x01}\\
 U&\mathtt{0x02}
\end{array}
\]
where $\operatorname{LE32}_{\rm bytes}$ encodes an integer in four
little-endian bytes.  This needs no secret setup ceremony.  It does not
remove the assumptions about the resulting points or their derivation.

For a column already stored as evaluations $p(\omega^i)$, precompute
Lagrange-basis generators
\[
 G_i^{\rm Lag}=\frac1n\sum_{j=0}^{n-1}[\omega^{-ij}]G_j.
\]
Here the factor $1/n$ also acts by scalar multiplication.  Inverse
Fourier interpolation gives
\[
 \Comm(p;r)=\sum_i[p(\omega^i)]G_i^{\rm Lag}+[r]W.
\]
Both bases therefore commit to the same polynomial.  Fixed and instance
commitments use the implementation's known default blind, which is
\emph{one}; it is not a fresh secret mask and not zero.

\subsection{The opening claim and its mask}

An opening claim consists of a commitment $P$, a point $x\in\F$ and a
value $v\in\F$.  The honest prover knows $(p,r)$ with
\[
 P=\Comm(p;r),\qquad p(x)=v,\qquad \deg p<n.
\]
The desired verification statement combines commitment consistency with
evaluation correctness.  Sending all coefficients would prove the
identity but disclose the polynomial and cost $n$ scalars.

The implemented opening starts by sampling a uniform degree-less-than-$n$
polynomial $s_0$ and setting
\[
 s(X)=s_0(X)-s_0(x).
\]
Thus $s(x)=0$.  This is a root at the opening point $x$, not necessarily
at zero.  The map $s_0\mapsto s$ is onto the kernel of evaluation at $x$,
with exactly $|\F|$ preimages per output, so $s$ is uniform in that
$(n-1)$-dimensional subspace.  The linear-map facts are proved in the
\emph{Math Guide}, ``Linear algebra and uniform masks''.

The prover samples an independent blind $r_s$ and sends
$S=\Comm(s;r_s)$.  The transcript then produces challenges $\xi$ and $z$,
in that order.  Define
\[
 \widehat p=p-v+\xi s,\qquad
 \widehat r=r+\xi r_s,\qquad
 \widehat P=P-[v]G_0+[\xi]S.
\]
The honest claim now says
$\widehat P=\Comm(\widehat p;\widehat r)$ and $\widehat p(x)=0$.
For a fixed nonzero $\xi$ independent of $s$, the polynomial
$p-v+\xi s$ is uniform in the zero-evaluation subspace, regardless of
which $p$ with evaluation $v$ was used.  This is a local masking fact,
not yet a simulator for an adaptively hashed transcript.

\subsection{One folding step}

Let $\boldsymbol a$ be the coefficients of $\widehat p$, and let
\[
 \boldsymbol b=(1,x,\ldots,x^{n-1}),\qquad f=\widehat r.
\]
Initially the verifier's group expression is $M=\widehat P$.
The invariant to preserve is
\begin{equation}\label{eq:ipa-invariant}
 M=\ip{\boldsymbol a}{\boldsymbol G}
   +[z\ip{\boldsymbol a}{\boldsymbol b}]U+[f]W.
\end{equation}
The inner product is initially zero, so this invariant holds for an
honest opening.

Split each vector into equal left and right halves, denoted $L,R$.
Sample fresh independent blinds $\lambda,\rho$, and send
\begin{align*}
 L_j&=\ip{\boldsymbol a_R}{\boldsymbol G_L}
       +[z\ip{\boldsymbol a_R}{\boldsymbol b_L}]U+[\lambda]W,\\
 R_j&=\ip{\boldsymbol a_L}{\boldsymbol G_R}
       +[z\ip{\boldsymbol a_L}{\boldsymbol b_R}]U+[\rho]W.
\end{align*}
The transcript supplies challenge $u_j$.  For $u_j\ne0$, fold by
\begin{align*}
 \boldsymbol a'&=\boldsymbol a_L+u_j^{-1}\boldsymbol a_R,&
 \boldsymbol b'&=\boldsymbol b_L+u_j\boldsymbol b_R,\\
 \boldsymbol G'&=\boldsymbol G_L+[u_j]\boldsymbol G_R,&
 f'&=f+u_j^{-1}\lambda+u_j\rho,\\
 M'&=M+[u_j^{-1}]L_j+[u_j]R_j.
\end{align*}
Vector operations are componentwise.

\begin{lemma}[The fold preserves the invariant]
If \eqref{eq:ipa-invariant} holds, then the same identity holds for the
primed vectors and scalars.
\end{lemma}
\begin{proof}
Expand the folded scalar--point inner product:
\[
\begin{split}
 \ip{\boldsymbol a'}{\boldsymbol G'}
 ={}&\ip{\boldsymbol a_L}{\boldsymbol G_L}
    +\ip{\boldsymbol a_R}{\boldsymbol G_R}\\
 &+[u_j]\ip{\boldsymbol a_L}{\boldsymbol G_R}
  +[u_j^{-1}]\ip{\boldsymbol a_R}{\boldsymbol G_L}.
\end{split}
\]
The scalar inner product has the identical expansion with
$\boldsymbol b$ in place of $\boldsymbol G$.  The two cross terms are
exactly those added by $[u_j^{-1}]L_j+[u_j]R_j$, including their blind
contributions.  The remaining terms are the old invariant.
\end{proof}

This asymmetric convention uses $u_j$ and $u_j^{-1}$, not their squares.
Changing to a symmetric convention without changing every equation
would produce a different protocol.

\subsection{The final check}

After $k=\log_2 n$ rounds, each vector has length one.  The prover sends
$c=a_{\rm final}$ and $f=f_{\rm final}$.  The verifier checks
\begin{equation}\label{eq:ipa-final}
 \widehat P+\sum_{j=0}^{k-1}
 \bigl([u_j^{-1}]L_j+[u_j]R_j\bigr)
 =
 [c]G_{\rm final}+[czb_{\rm final}]U+[f]W.
\end{equation}
It calculates the folded public quantities without receiving them.
Define the degree-$n-1$ polynomial
\[
 B(X)=\prod_{j=0}^{k-1}
       \left(1+u_jX^{\,n/2^{j+1}}\right)
       =\sum_{i=0}^{n-1}B_iX^i.
\]
Then
\[
 b_{\rm final}=B(x),\qquad
 G_{\rm final}=\sum_{i=0}^{n-1}[B_i]G_i.
\]
The first split separates coefficient indices by $n/2$, the next by
$n/4$, and so on, which gives this factor order.  For $n=4$, for example,
the coefficient weights are $(1,u_1,u_0,u_0u_1)$.

The opening proof contains $S$, two group elements per round, and the
two final scalars.  Its communication is logarithmic in $n$, but the
ordinary verifier still computes a linear-size multiscalar
multiplication to obtain the final generator contribution.  A
multiscalar multiplication is a sum of specified scalar multiples of
points; batching such sums does not itself verify a proof inside another
proof.

The implementation returns a deferred verification object called a
\texttt{Guard}.  Its challenge-dependent generator contribution must
still be supplied and the resulting multiscalar multiplication evaluated
to the identity.  The single-proof verification strategy does both.
Constructing the deferred object alone is not proof acceptance.

The algebra above excludes zero challenges where an inverse is needed;
a zero $z$ or $\xi$ also removes a needed security contribution.  These
are exceptional events in the analysis.  The pinned transcript does
not implement an implicit retry-until-nonzero convention.  An
implementation failure or rejection on an exceptional transcript must
not be described as resampling.

\section{Reducing multiple openings}\label{sec:multiopen}

\subsection{The list to be authenticated}

The gate, permutation and lookup checks query many polynomials at
several rotations of $x$.  The verifier has an ordered list
\[
 (P_i,t_i,v_i),
\]
where each $P_i$ names a particular committed polynomial, $t_i$ is an
evaluation point, and $v_i$ is the claimed value.  Repeated references to
one polynomial use the same identity.  Each polynomial--point pair
occurs only once in this input list.  The verification key and the
argument construction determine that list, not an unauthenticated
prover-supplied description.

Different polynomials can have different point sets.  For example, an
advice column might be queried only at $x$, while a permutation product
is queried at $x,\omega x,\omega^\ell x$.  Opening each query separately
would repeat the inner-product argument.  Halo~2 reduces all the queries
to one opening using polynomial interpolation and linear commitments.

\subsection{First combine equal point sets}

Collect the distinct points at which each polynomial is queried.
Group polynomials with identical point sets into groups
$j=0,\ldots,m-1$.  The implementation preserves first occurrence of each
polynomial in the query list; groups are ordered by the first polynomial
belonging to each group.  Point order follows first occurrence in the
original list.  These rules make both sides derive the same order.

After the original evaluations have entered the transcript, derive
$x_1$ and $x_2$.  If group $j$ contains polynomials
$p_{j,0},\ldots,p_{j,k_j-1}$ in that order, define
\[
 Q_j(X)=\sum_{i=0}^{k_j-1}x_1^{k_j-1-i}p_{j,i}(X).
\]
Use the same Horner combination for their commitments, blinds and
claimed values at every point of the shared set $T_j$.  For instance,
three polynomials combine as $x_1^2p_0+x_1p_1+p_2$, not
$p_0+x_1p_1+x_1^2p_2$.

Let $r_j$ be the unique interpolation polynomial of degree less than
$|T_j|$ agreeing with these combined claimed values, and let
\[
 V_j(X)=\prod_{t\in T_j}(X-t).
\]
The entire grouped claim is equivalent to
\[
 V_j(X)\mid Q_j(X)-r_j(X).
\]
The \emph{Math Guide}, ``Polynomials and evaluation'', constructs both
interpolation and division, so this is an explicit polynomial test,
not an assumed new commitment primitive.

\subsection{Combine the divisibility claims}

For honest claims, set
\[
 f_j(X)=\frac{Q_j(X)-r_j(X)}{V_j(X)},\qquad
 Q'(X)=\sum_{j=0}^{m-1}x_2^{m-1-j}f_j(X).
\]
Every quotient has degree less than $n$.  The prover commits to $Q'$
with a fresh independent blind and sends this one commitment.

An implementation can compute $f_j$ by repeated synthetic division of
$Q_j$ by the factors $X-t$, discarding remainders.  The resulting
quotient is also the quotient of $Q_j-r_j$: the unique remainder upon
division by $V_j$ is precisely the interpolation polynomial of $Q_j$'s
actual values.  If a claimed value is wrong, the prover's actual
remainder and the verifier's claimed interpolation differ.  The next
check is where that discrepancy must be caught.

After the commitment to $Q'$, derive $x_3$.  The prover sends
\[
 u_j=Q_j(x_3)\qquad(0\le j<m).
\]
Provided $x_3\notin\bigcup_jT_j$, the verifier computes the expected value
\begin{equation}\label{eq:multiopen-expected}
 v'=\sum_{j=0}^{m-1}x_2^{m-1-j}
       \frac{u_j-r_j(x_3)}{V_j(x_3)}.
\end{equation}
The claim is that $Q'(x_3)=v'$.  The challenge $x_3$ is sampled and the
claims $u_j$ are transmitted only after the commitment to $Q'$.  Some
values can already be predictable, for example for a constant $Q_j$;
the required order concerns the challenge and messages.

For a singleton set $T_j=\{t\}$, interpolation is just the constant
claimed value $r_j=v_j$, and the summand is
$(u_j-v_j)/(x_3-t)$.  For two points $t_0,t_1$, it is obtained using
\[
 r_j(X)=v_{j,0}\frac{X-t_1}{t_0-t_1}
       +v_{j,1}\frac{X-t_0}{t_1-t_0},
 \qquad V_j(X)=(X-t_0)(X-t_1).
\]
Thus both the single-point and rotated-query cases use the same
construction.

\subsection{One final polynomial}

Only after the $u_j$ values are in the transcript, derive $x_4$.
Starting with $P_{\rm final}=Q'$, perform the Horner update
$P_{\rm final}\leftarrow x_4P_{\rm final}+Q_j$ for groups in order.
Equivalently,
\[
 P_{\rm final}(X)=x_4^mQ'(X)
       +\sum_{j=0}^{m-1}x_4^{m-1-j}Q_j(X).
\]
Its commitment is the same linear combination of known commitments;
no extra final commitment is sent.  The verifier calculates the claimed
evaluation
\[
 v_{\rm final}=x_4^mv'
       +\sum_{j=0}^{m-1}x_4^{m-1-j}u_j.
\]
The prover now supplies the masked inner-product opening of
$P_{\rm final}$ at $x_3$ to $v_{\rm final}$, exactly as in
\S\ref{sec:ipa}.  Linear combination of commitment blinds gives the
prover's final opening blind.

Each challenge has a specific job.  The challenge $x_1$ prevents fixed
errors among polynomials in a group from cancelling deliberately;
$x_2$ combines the quotient identities; $x_3$ tests their agreement at
a fresh point; $x_4$ combines the last evaluation claims.  These
explanations rely on earlier objects being fixed and on the eventual
opening argument having the necessary soundness guarantee.  They do
not turn computational binding by itself into knowledge extraction.

Point coincidences, including $x_3\in T_j$, belong to the exceptional
set.  Their actual implementation handling is not a licence to divide
by zero in the proof of correctness.  Nor may an analysis assume that
every differently named rotation is distinct when $x=0$.

\section{The transcript and the verifier}\label{sec:transcript}

\subsection{The circuit-specific keys}

The public parameters depend on the field, group and domain size.  A
verification key additionally fixes the circuit: column counts, query
rotations, gate and lookup expressions, equality wiring, fixed
commitments and degree.  A proving key contains that verification key
and precomputed polynomial forms needed by the prover.  Key generation
does not require the private transaction witness.  The resulting
circuit shape must not change with the witness.

The verifier must obtain the intended verification key through its
application's authenticated rules.  Accepting a prover-selected key
would allow a proof for a different, weaker relation.  Fixing the key
also fixes the meaning and position of every public input.

In the pinned implementation the verification-key transcript
representative is formed by taking the Rust debug-format string of the
key's \emph{pinned} contents, prefixing its UTF-8 byte length as an
eight-byte little-endian integer, hashing with 64-byte BLAKE2b under
personalisation \texttt{Halo2-Verify-Key}, and reducing the 64-byte
little-endian result modulo the scalar-field order.  The pinned
contents include both field moduli, domain parameters, constraint
system, fixed commitments and permutation key.  This is a
version-specific circuit identifier, not a claim that arbitrary Rust
debug formatting is a portable consensus serialisation.

The verifier constructs each public instance polynomial by zero-padding
the supplied column to $n$ entries, with at most $\ell$ supplied entries.
It computes its commitment with known blind one.  The verification-key
representative and these instance commitments enter the transcript as
common inputs, without being transmitted as proof bytes.

\subsection{Hashing messages into challenges}

The \emph{Crypto Guide}, ``Hash functions and key derivation'', constructs
BLAKE2b.  Halo~2 initialises its transcript with 64-byte BLAKE2b output
and personalisation \texttt{Halo2-Transcript}, padded to the hash's
personalisation-field length.  The running hash state absorbs the
following typed messages:
\[
\begin{array}{c|l}
 \text{message}&\text{bytes absorbed}\\ \hline
 \text{scalar }a&\mathtt{0x02}\parallel a^\star\\
 \text{point }P&\mathtt{0x01}\parallel x(P)^\star
                       \parallel y(P)^\star\\
 \text{challenge request}&\mathtt{0x00}
\end{array}
\]
Scalar and coordinate encodings here are canonical 32-byte
little-endian field encodings.  On a challenge request, the state first
absorbs the zero byte, then a clone of the state is finalised.  Reduce
the resulting 64-byte little-endian integer modulo $|\F|$ to obtain the
challenge.  The original state retains the newly absorbed zero byte;
the hash output itself is not appended to it.

Proof bytes and hashed bytes are not identical for points.  The proof
contains the canonical 32-byte compressed point, whereas the transcript
absorbs its two affine coordinates.  Decoding must succeed before the
coordinates can be hashed.  The transcript rejects the identity point,
which has no affine coordinates, even though the curve encoding format
can represent that identity.  A scalar in the proof must also decode
canonically; reducing an out-of-range proof encoding would be a
different parser.

The point and scalar prefix bytes separate their meanings in the hash
input.  Circuit and instance binding comes from the common-input prefix.
Challenge ordering comes from the prescribed message sequence, not from
serialising the names $\theta,\beta$ and so on.

Reduction of a uniformly random 512-bit integer is negligibly biased
rather than exactly uniform unless the field order divides $2^{512}$.
The local algebraic probability calculations use ideal uniform
challenges.  A concrete security estimate includes this statistical
difference as well as the hash model and the number of adversarial
hash queries.

\subsection{The whole sequence}

Before combining constraints with $y$, the prover additionally samples
a uniform degree-less-than-$n$ polynomial $R$ and an independent
commitment blind, and sends $C_R=\Comm(R;r_R)$.  It later sends $R(x)$.
Both $R$ and the collapsed quotient $h_x$ are included among the
polynomials opened at $x$.  This masks information in the later
multi-opening reduction.  The polynomial $R$ is not added to the gate
numerator or to the quotient identity.

The transcript order is:
\begin{center}
\begin{tabular}{p{0.64\linewidth}|p{0.20\linewidth}}
 Common inputs or prover messages&Next challenges\\ \hline
 Verification-key representative; instance commitments&None\\
 Advice commitments&$\theta$\\
 Permuted lookup commitments $(A',S')$&
 $\beta,\gamma$\\
 Permutation products; lookup products; $C_R$&$y$\\
 Quotient chunk commitments&$x$\\
 Column and auxiliary evaluations&$x_1,x_2$\\
 Combined quotient commitment $C_{Q'}$&$x_3$\\
 Group evaluations $u_0,\ldots,u_{m-1}$&$x_4$\\
 Opening mask commitment $S$&$\xi,z$\\
 Folding pair $L_j,R_j$, once per round&$u_j$\\
 Final scalars $c,f$&Final group check
\end{tabular}
\end{center}
For several circuit instances proved together, each stage includes all
instances before proceeding to the next challenge.  Fixed columns,
fixed permutation commitments, the random polynomial and the aggregate
quotient are shared where the construction specifies; advice and
auxiliary products remain per-instance.  This shares one argument
over several witnesses, not a proof that verifies earlier proofs.

The scalar evaluation messages after $x$ are ordered as follows:
\begin{enumerate}[nosep]
 \item Instance-column queries, in configured order for each instance.
 \item Advice-column queries, likewise.
 \item Shared fixed-column queries.
 \item The scalar $R(x)$.
 \item The fixed permutation-polynomial values $s_j(x)$.
 \item For each instance and each permutation chunk, the values at
       $x,\omega x$, then at $\omega^\ell x$ unless it is the last chunk.
 \item For each instance and lookup, the five values
       $Z(x),Z(\omega x),A'(x),A'(\omega^{-1}x),S'(x)$.
\end{enumerate}
There is no separate message for the expected quotient value
$N_y(x)/(x^n-1)$.

The multi-opening query list has its own deterministic order.  For
each instance it lists instance queries, advice queries, permutation
queries, then lookup queries; it then appends the shared fixed queries,
fixed permutation queries, and finally $h_x$ and $R$ at $x$.
Permutation queries first list $(x,\omega x)$ for each chunk, then the
$\omega^\ell x$ queries for the preceding chunks in reverse order.
Each lookup lists $Z,A',S'$ at $x$, then $A'$ at $\omega^{-1}x$, then
$Z$ at $\omega x$.  This order determines the point-set grouping and
Horner weights in \S\ref{sec:multiopen}; it need not equal the scalar
transmission order.

\subsection{The acceptance algorithm}

Given the authenticated verification key, parameters, public instance
columns and a proof byte slice, the verifier performs the following
logical steps.
\begin{enumerate}
 \item Check instance column counts and lengths.  Construct the common
       inputs and initialise the transcript.
 \item Read canonical proof elements in the fixed sequence, rejecting
       malformed encodings and read failures.  Derive each challenge
       only after the preceding messages have been absorbed.
 \item Evaluate every gate expression at $x$ using the claimed column
       evaluations.  Evaluate every permutation and lookup constraint
       using its claimed auxiliary evaluations and the public
       $L_0(x),L_\ell(x),L_{\rm blind}(x)$.
 \item Combine those expressions in Horner order with $y$.  Compute
       the expected quotient evaluation and the commitment to $h_x$.
 \item Construct the complete opening-query list.  Reproduce the
       point-set groups, interpolants and linear commitment
       combinations.  Compute the expected final evaluation at $x_3$.
 \item Read and verify the masked inner-product opening.  Supply the
       final challenge-dependent generator contribution and evaluate
       the entire multiscalar multiplication to $\mathcal O$.
\end{enumerate}
The surrounding application supplies proof framing and enforces its
specified length or end-of-input rule.  The generic transcript reader
consumes the elements it needs; it does not by itself establish that
an arbitrary enclosing byte stream has no trailing data.

Evaluations are claims until the final opening check succeeds.  A
vanishing test performed on unauthenticated scalar messages is not
already a verified circuit.  Conversely, an opening proof authenticates
a polynomial claim, not the application's intended semantics: the
circuit and public-input binding remain essential.

\section{Knowledge soundness and zero knowledge}\label{sec:security}

\subsection{What the algebra has established}

The preceding derivations establish deterministic identities:
interpolation represents columns, vanishing means divisibility, the
permutation and lookup constraints enforce their advertised relations
outside specified bad events, and every honest inner-product fold
preserves the group equation.  They therefore explain honest
correctness and the structure of the verification test.

A malicious prover need not generate group elements from known
coefficient vectors using the honest algorithm.  For example, merely
observing a point $P$ does not reveal a polynomial and blind opening it.
The expression $\Comm(p;r)=P$ has many mathematical solutions.
An \emph{extractor} is an algorithm with specified access to the prover
which obtains a suitable witness efficiently; saying that a solution
exists is not a construction of such an algorithm.

This distinction is why the folding invariant is not, by itself, a
knowledge-soundness proof.  It checks what happens if the vectors used
in the calculation exist with the requisite consistency.  Extraction
has to establish that consistency against an adversarial transcript.

\subsection{A conditional arithmetic-layer statement}

It is useful to isolate exactly what a secure commitment-and-opening
layer must provide.  Consider first an interactive version with
independent uniform verifier challenges.

\begin{assumption}[Consistent polynomial extraction]
Except with probability $\varepsilon_{\rm open}$, an extractor
associates every committed polynomial with coefficients of the claimed
degree, consistently across all its uses, and accepted openings have
the corresponding evaluations.  The extraction guarantee respects the
challenge chronology: it supplies the relevant fixed polynomial objects
before the later challenges to which the algebraic tests are applied.
\end{assumption}

This is a security requirement on the composed opening protocol, not a
property proved by the Pedersen commitment's hiding or binding lemma.
It also excludes interpreting a polynomial fitted after learning the
test point as a polynomial fixed before that point.

\begin{proposition}[Recovering a satisfying arithmetic witness]
Under consistent polynomial extraction, and outside polynomial-test,
compression and zero-factor bad events, an accepted interactive proof
yields advice values satisfying the circuit gates, copy equalities and
lookup memberships.
\end{proposition}
\begin{proof}
Evaluate the extracted advice polynomials on the usable rows of $H$.
Authentication of all queried evaluations makes the verifier's
computed numerator the value of the actual combined constraint
polynomial.  Outside the random-point identity-test event,
$N_y=(X^n-1)h$ as polynomials.  Outside the random-combination event,
each constituent constraint therefore vanishes on $H$.

The permutation recurrences telescope.  Their initial value is one,
their chunk links connect the products, and nonzero factors make the
Boolean final value equal one.  Outside the labelled-product
collision event, the permutation lemma implies all copy equalities.

The lookup products likewise prove that $A'$ and $S'$ have the same
multisets as the compressed input and table.  The first-row match and
run constraints then place every compressed input in the compressed
table.  Outside tuple-compression collisions, that is membership of the
original tuples.  The custom gates already vanish at every required
row.  The extracted advice values are the required circuit witness.
\end{proof}

All algebraic error terms can be bounded using the elementary root
bound.  For example, for $M$ combined constraints the $y$ test contributes
at most $(M-1)/|\F|$.  A committed degree-less-than-$(d-1)n$ quotient
makes the final difference have degree at most $dn-1$, giving at most
$(dn-1)/|\F|$ for the $x$ test.  With $C$ participating copy cells, a
nonzero labelled-product difference has total degree at most $C$ in
$\beta,\gamma$.  Each zero factor has probability at most $1/|\F|$;
count all numerator and denominator factors.  Lookup compression must
count all distinct input--table tuple pairs, not merely one pair.
A union bound combines these contributions and $\varepsilon_{\rm open}$.
Excluded challenges such as $x\in H$ must be included in the failure
accounting as well.

These are deliberately modular bounds, not an asserted concrete
security level for an entire production proof.  The application's
circuit-to-relation implication is a further step after extracting a
satisfying circuit witness.

\subsection{The commitment and Fiat--Shamir boundary}

The usual analysis of this family of inner-product commitments uses
the discrete-logarithm assumption and a model controlling adversarial
group computations.  In the \emph{algebraic group model}, an adversary
that outputs a group element must also supply, to the analysis, a
linear representation of it in previously available group elements.
The verifier does not receive these representations.  They are an
extra restriction on the analysed adversary, not a consequence of
canonical point decoding or of the group being cyclic.

Such representations let an extractor compare generator coefficients
in suitable related accepting transcripts.  Unknown efficient
relations between independently generated points prevent inconsistent
representations except with the model's failure probability.
Constructing the extractor, its rewinding strategy and its time bound
is a separate theorem obligation.  The \emph{Crypto Guide},
``Knowledge, simulation, and Schnorr'', demonstrates these distinctions
for a simpler three-message protocol; that demonstration is not
automatically a theorem for this longer protocol.

Fiat--Shamir replaces interactive challenges by transcript hashes.
Its analysis models the hash as an oracle with consistent random
answers on new inputs and counts the adversary's queries.  The
adversary can try different commitment prefixes and observe their
challenges.  Consequently, an interactive one-attempt root bound
cannot simply be relabelled as the noninteractive security bound.
Rewinding or oracle programming, when used by a proof, must be
justified within its exact model, including adaptively selected public
inputs and correlated openings.

The Halo paper supplies the originating inner-product construction
and its modelled security analysis.  The pinned implementation includes
a masked variant and the concrete reductions described here.  We do
not assert that a short folding identity or an informal appeal to
that paper proves every detail of this implementation's complete
composition.  The explicit extraction assumption above states the
remaining boundary instead of presenting it as a proved consequence
of binding.  This boundary matters when using the argument as a
security premise for the application.

\subsection{Why polynomial evaluations need masks}

A hiding commitment does not hide a subsequently disclosed evaluation.
If a witness polynomial were determined entirely by its active rows,
each evaluation could reveal a linear equation in those rows.  Halo~2
therefore randomises unused row values as well as commitment blinds.

Fix a set $B$ of $s$ independently randomised domain points and fix all
other row values.  Let
\[
 T_B(X)=\prod_{h\in H\setminus B}(X-h).
\]
The possible masked column polynomials form the affine space
\[
 a_0(X)+T_B(X)r(X),\qquad \deg r<s.
\]
Choosing the $s$ row masks uniformly chooses $r$ uniformly: both
parameterisations describe the same $s$-dimensional affine space
bijectively.

For $t\le s$ fixed distinct query points $z_1,\ldots,z_t$ outside $H$,
the evaluation vector of the mask is
\[
 \bigl(T_B(z_1)r(z_1),\ldots,T_B(z_t)r(z_t)\bigr).
\]
Interpolation makes $r\mapsto(r(z_1),\ldots,r(z_t))$ surjective, and all
the factors $T_B(z_i)$ are nonzero.  The \emph{Math Guide}, ``Linear
algebra and uniform masks'', therefore implies that this vector is
uniform on $\F^t$.  It hides the corresponding evaluations of $a_0$
jointly, for these fixed points and this isolated linear observation.

The implementation chooses
\[
 b=\max(3,q_{\rm adv})+2,
\]
where $q_{\rm adv}$ is the maximum number of distinct configured
rotations of any advice column, taking one as the fallback when there
are no advice columns.  The minimum three covers permutation
polynomials queried at three rotations.  One additional factor covers
the multi-opening point $x_3$, and one is an extra margin.  Advice and
permuted lookup columns actually randomise $b+1$ rows; product
polynomials randomise $b$.  This accounts for the concrete row budget
rather than merely saying that unspecified blinding is added.

\subsection{The quotient mask and the joint transcript}

The quotient depends nonlinearly on witness columns and auxiliary
products.  Its evaluation $h(x)$ is already determined by the
constraint evaluations, but an additional value at $x_3$ could disclose
new information.  The independent random polynomial $R$ is included
with $h_x$ in the group of polynomials queried at $x$.  Their later
Horner-combined evaluation at $x_3$ contains a random contribution from
$R(x_3)$.

For fixed distinct $x,x_3$ and $n\ge2$, the pair
$(R(x),R(x_3))$ is uniform in $\F^2$.  After revealing the first
coordinate, the second remains uniform.  When its batching coefficient
is nonzero, it can therefore mask another scalar in that linear
combination.  The PCS mask $s$ separately masks the final opened
polynomial while preserving its required evaluation.

Neither observation proves full zero knowledge in isolation.
Different messages can share masks, challenges are derived from
earlier commitments, and later openings must remain consistent with
those commitments.  Two individually uniform messages can jointly
reveal a witness, as the example $(r,r+w)$ shows.  A complete
zero-knowledge proof must construct one joint simulator for the
entire view, including commitments, all evaluations, quotient
reductions, inner-product messages and hash-derived challenges.
If it relies on programming the hash oracle, that ability belongs in
the statement of the theorem.

The local affine-mask and folding lemmas explain the mechanisms and
their required dimensions.  This volume does not replace the joint
simulation obligation by those lemmas, nor does it claim that any
missing composition theorem has thereby been proved.

\subsection{The application-facing contract}

To consume Halo~2 correctly, an application must fix the circuit and
verification key, define the public-input mapping, frame and parse the
proof, run the complete opening verification, and check every
non-circuit rule separately.  It must establish that a satisfying
circuit witness implies its intended relation, with canonical
encodings and integer bounds where needed.

The resulting security claim is conditional on that circuit argument,
the commitment and transcript security assumptions, correct parameter
and key derivation, and correct implementation.  Proof privacy does
not hide the public input, the proof's circuit shape, transaction
metadata outside the circuit, or information the wallet discloses
through its network activity.  These are boundaries of the proved
relation, not failures of the masking algebra.

\subsection{Primary construction sources}\label{sec:halo-sources}

The mechanism and byte order in this volume are checked against
\href{https://github.com/zcash/halo2/tree/261faaccd5a30c19bb8468600841a0dac305adf5/halo2_proofs/src}
{the pinned \texttt{halo2\_proofs} source}.
Within that source, the relevant files are:
\begin{itemize}[nosep]
 \item \texttt{plonk/circuit.rs}, \texttt{plonk/prover.rs} and
       \texttt{plonk/verifier.rs}: circuit degree, blinding budget,
       complete message and query order.
 \item \texttt{plonk/permutation/} and \texttt{plonk/lookup/}:
       product recurrences, exact Boolean boundaries and rotations.
 \item \texttt{poly/domain.rs} and \texttt{plonk/vanishing/}:
       extended coset, random polynomial and quotient chunks.
 \item \texttt{poly/commitment.rs}, \texttt{poly/commitment/} and
       \texttt{poly/multiopen/}: generator derivation, asymmetric
       masked folding and Horner-ordered multi-opening reduction.
 \item \texttt{transcript.rs} and \texttt{plonk.rs}: challenge bytes
       and verification-key representative.
\end{itemize}
The
\href{https://zcash.github.io/halo2/design/proving-system.html}
{Halo~2 design documentation} explains the design; its sketches are not
substituted for the concrete source where conventions differ.
Bowe, Grigg and Hopwood's
\href{https://eprint.iacr.org/2019/1021}
{\emph{Halo: Recursive Proof Composition without a Trusted Setup}}
is the primary paper behind the commitment construction and modelled
security analysis.  Recursive composition itself is not needed to
understand the single proof verified here.

The executable checks in \texttt{research/check\_foundations.py} and
\texttt{research/check\_halo2.py} recompute the small-field example,
folding identities and multi-opening algebra.  They test the displayed
calculations, not cryptographic hardness, witness extraction or a
complete transcript simulation.

\end{document}
